% Generated draft: requires mathematical and source review.
\documentclass[11pt]{article}
\usepackage[margin=25mm]{geometry}
\usepackage{fontspec}
\setmainfont{Latin Modern Roman}
\usepackage{amsmath,amssymb,mathtools,unicode-math}
\setmathfont{Latin Modern Math}
\usepackage{enumitem,needspace,xurl,hyperref}
\setlength{\parindent}{0pt}
\setlength{\parskip}{5pt}
\setlength{\emergencystretch}{3em}
\newcommand{\N}{\mathbb{N}}
\newcommand{\Z}{\mathbb{Z}}
\newcommand{\Q}{\mathbb{Q}}
\newcommand{\R}{\mathbb{R}}
\newcommand{\C}{\mathbb{C}}
\newcommand{\F}{\mathbb{F}}
\newcommand{\A}{\mathbb{A}}
\newcommand{\PP}{\mathbb{P}}
\newcommand{\E}{\mathbb{E}}
\newcommand{\eps}{\varepsilon}
\newcommand{\dd}{\,\mathrm{d}}
\begin{document}
\section*{UnsolvedMath ID: 20002716}
\textbf{Source problem code:} catalogue entry
\subsection{Definitions and mathematical statement}
P. Dehornoy's Problem 4.2, split typographically as ``Problem 4. 2.,'' has three parts. (1) Is a cancellative finitely generated monoid with lcm's necessarily Garside, equivalently does it have a Garside element? (2) Can the Coxeter-group/Cayley-graph source of good Artin Garside structures be systematized for general Garside groups? (3) For an atomic Garside monoid, a Garside element $\Delta$ and a length $\ell$, is
\[
\ell(\Delta^k)\leq Ck|\Delta|\qquad(k\geq1)
\]
true for a constant independent of $k$? The source omits whether both left and right lcm's are required, unit/Noetherian hypotheses, the kind of Coxeter-like object, and the meanings of $\ell$ and $|\Delta|$.

For the first part the aid supplies a sufficient atom-level certificate. Work with a unit-free atomic cancellative monoid generated by a finite atom set $\mathcal A$, with the relevant left and right lcm hypotheses. An atom is a nonunit not decomposable into two nonunits. Let $\Delta$ be the common right lcm of the atoms, so each atom left-divides $\Delta$. If an atom permutation $\tau$ satisfies
\[
a\Delta=\Delta\tau(a)\qquad(a\in\mathcal A),
\]
then the aid asserts that $\Delta$ is a Garside element: its left and right divisor sets agree, are finite, and generate the monoid. These balanced-divisor and finiteness conditions are part of the Garside convention. The certificate adds quasi-centrality; it does not show that the original lcm assumptions always imply it.

For the second part the aid rules out a specific universal recipe, not every Coxeter analogue. If a marked Garside structure has $\Delta=a^2$ for an atom $a$, any quotient making all atoms involutions maps both $1$ and $\Delta$ to the identity. Hence it cannot be injective on $\operatorname{Div}(\Delta)$. The example
\[
M=\langle a,b\mid a^2=b^2\rangle^+,\qquad\Delta=a^2=b^2,
\]
has the standard Garside divisor set $\{1,a,b,\Delta\}$ and gives this obstruction to preserving the simple-divisor geometry via the naive atom-square quotient. It does not prove absence of all possible associated Coxeter-like objects.

For the third part the aid distinguishes three length conventions. If $\ell$ is shortest atomic word length and $|\Delta|=\ell(\Delta)$, subadditivity gives the claim with $C=1$. If instead ``length'' permits an arbitrary superadditive atomicity norm, $M=\langle a\rangle^+$, $\Delta=a$ and $\ell(a^j)=j^2$ give a quadratic counterexample to a uniform linear bound; this is not shortest or maximal factorization length. Finally let $\ell_{\max}(x)$ be the maximal length of an atomic factorization and suppose a positive additive grading $g:M\to\N_0$ exists, with $g(a)>0$ for all atoms. With $m=\min_{a\in\mathcal A}g(a)>0$,
\[
\ell_{\max}(\Delta^k)\leq\frac{k\,g(\Delta)}m.
\]
This yields the requested linear form for any fixed positive reference $|\Delta|$, with the corresponding constant. It does not settle maximal atomic length in the absence of a positive grading. The undefined source length convention is therefore decisive for interpreting a positive or negative answer.
\subsection{Short English statement}
Dehornoy asks about Garside existence, Coxeter-like models, and linear growth of lengths of Garside powers. The aid gives a quasi-centrality certificate, an obstruction to the naive involution quotient, and different answers for different length notions. None of these distinctions should be collapsed into one general solution.
\subsection{Sources}
\begin{itemize}
\item Local UnsolvedMath registry, permanent ID $20002716$, \texttt{catalogue entry}; \texttt{aim-probability-notes.json}, \texttt{notes[157]}, extracted source item 4.
\item American Institute of Mathematics, \emph{Braid groups, clusters and free probability}, Problem 4.2 (P. Dehornoy), \url{https://aimath.org/WWN/braidgroups/braidgroups.pdf}. Supplied verification repairs the split heading and the line-break ``neces-sarily.''
\item Research attempt 1: partial result, problem partially solved at run, rigor label ``proved,'' novelty confidence low. Registry classification partial progress, status partially solved, statement status exact. No authoritative separate references are supplied for the atom certificate, quotient example or length results.
\end{itemize}
\end{document}
